% 《大学物理学学习指导（第5版）》第2章：例2.1。
% 源 PDF 第28页，书内印刷页20。质量状态：待独立复核。

\sourceexample{2.1}{20}{28}
\problemtext 如原书例2.1的轻绳和滑轮系统：质量分别为$m_1,m_2$的物块$A,B$放在无摩擦水平桌面上，
一根轻绳的两端连接$A,B$，中部绕过理想轻动滑轮；另一轻绳悬挂滑轮并与竖直悬物$C$（质量$m_3$）相连。
忽略绳与滑轮质量和摩擦，绳不可伸长，求三物块加速度大小。
\workings
设动滑轮（及$C$）向下加速度$a$；$A,B$沿桌面向滑轮的加速度大小分别为$a_A,a_B$。
绳长约束使
\[
a_A+a_B=2a.
\]
设桌面绳张力为$T'$，悬物绳张力为$T$。轻滑轮受力平衡给出$T=2T'$。
沿各自运动方向应用牛顿第二定律，
\[
m_1a_A=T',\quad m_2a_B=T',\quad m_3g-T=m_3a.
\]
于是$a_A=T'/m_1,\ a_B=T'/m_2$，由绳长约束得
\[
2a=T'\left(\frac1{m_1}+\frac1{m_2}\right)
\quad\Longrightarrow\quad
T'=\frac{2am_1m_2}{m_1+m_2}.
\]
代入$m_3g-2T'=m_3a$，
\[
\boxed{
a=\frac{m_3(m_1+m_2)g}{4m_1m_2+m_3(m_1+m_2)}
}.
\]
进而
\[
\boxed{
a_A=\frac{2m_2m_3g}{4m_1m_2+m_3(m_1+m_2)},\qquad
a_B=\frac{2m_1m_3g}{4m_1m_2+m_3(m_1+m_2)}.
}
\]
\crosscheck (i) 各式量纲均为加速度；(ii) $a_A+a_B=2a$；
(iii) 若$m_1=m_2$，则$a_A=a_B=a$；
(iv) $m_3\to0$时三个加速度均趋近于零。
这是原书图意的文字表达，正式制图校对仍应逐项核验绳的绕行方式。
